\section{Introduction}\label{sec:intro}

Artificial intelligence (AI) is increasingly presented as a tool for climate adaptation. Machine
learning now supports early warning of floods, droughts and heatwaves, crop-yield forecasting,
pest detection and the targeting of relief after disasters
\citep{rolnick2022,lealfilho2022,reichstein2025}. Development agencies, multilateral banks and
governments are beginning to fund ``AI for climate resilience'' programmes in low- and
middle-income countries. Every such programme faces the same allocation problem: capacity cannot
be built everywhere at once, so \emph{where should it be built first?}

The answer depends on whose resilience counts. Climate damages fall disproportionately on poor
countries and on poor households within them \citep{hallegatte2017,diffenbaugh2019}, whereas the
capabilities needed to develop and use AI are concentrated in rich countries and better-off groups
\citep{cazzaniga2024,korinek2021}. If AI-enabled adaptation protects mainly the incomes of
households that are already well connected and well off, a programme that maximises total
resilience may leave the most vulnerable no better protected \citep{larosa2024}. AI capability
also has a physical cost: data-centre electricity use is projected to more than double by 2030
\citep{iea2025,devries2023}, and the carbon-abatement benefits of AI appear weaker where income is
more unequal \citep{dong2024}. The special issue asks who benefits from AI-driven climate solutions,
who bears their costs, and under what conditions AI promotes inclusive climate-resilient
development. These are questions about distribution and about targeting, and the standard
toolkit of the empirical literature, which estimates an average effect and adds an interaction
term with a single inequality measure, is poorly suited to answer them.

This paper develops and applies a method designed for the targeting question. We call it
\emph{equity-weighted doubly robust policy learning} (EW-DRPL). It has three parts. First, it uses
year-to-year temperature shocks, which are plausibly exogenous once country-specific trends are
removed \citep{dell2012,burke2015}, to measure how sensitive the income growth of each decile of a
country's distribution is to climate shocks. Second, it estimates how this sensitivity differs
between countries with and without climate-AI capability through orthogonalised, cross-fitted,
doubly robust scores \citep{robinson1988,chernozhukov2018}. These scores let machine-learning
methods absorb the strong confounding between AI capability and development without biasing the
estimated resilience effect. Third, it aggregates decile-level resilience effects with
Atkinson-type welfare weights \citep{atkinson1970} and learns an interpretable allocation rule, a
shallow decision tree, that maximises equity-weighted welfare under a budget constraint
\citep{kitagawa2018,athey2021}. Varying the inequality-aversion parameter shows how the
recommended allocation changes as a planner cares more about the poorest. The difference between
the rules yields two quantities that speak directly to the special issue: the \emph{equity gain} of
targeting for the poor, and the \emph{equity price} in aggregate resilience of doing so.

We apply EW-DRPL to a new global dataset, focusing on agriculture, the sector through which heat
most directly reaches poor households \citep{ortizbobea2021,christiaensen2011}. National
climate-AI capability is measured from the OpenAlex index of scholarly works \citep{priem2022} as a
depreciated stock of research at the intersection of AI or machine learning and climate change.
A country counts as capable when its stock exceeds the level predicted by its income, so that
capability is not simply a proxy for development. We also measure \emph{knowledge sovereignty}, the
share of a country's climate-AI research carried out by domestic institutions alone. Outcomes are
cereal-yield growth for 145 countries from 2000 to 2023, shocks are hottest-month temperature
anomalies from the ERA5 reanalysis, and the welfare weights use decile incomes from the World
Bank's Poverty and Inequality Platform together with the agricultural exposure of each decile. A
Monte Carlo study shows how the estimator behaves at the sample sizes typical of cross-country
work, and an appendix applies the method directly to decile income growth from 1,334
household-survey spells.

The results are sobering. A 1$^{\circ}$C hotter-than-trend hottest month lowers cereal-yield growth
by 3.6 percentage points, but countries with climate-AI capability suffer the same losses as
countries without it. The average resilience effect is 0.26 points (SE 1.07), and effects larger
than about two-thirds of the damage are ruled out. A standard interaction regression on
agricultural value added would suggest a marginally significant resilience dividend, which the
doubly robust estimates do not confirm. The allocation rule closest to current practice, building
on existing AI strength, delivers essentially no equity-weighted benefit, because the leading
countries have little agricultural exposure. Rules that favour the poorest and most agricultural
countries have higher, though imprecise, value. Inequality aversion changes the recommended
allocation only when it is high; at that point Viet Nam and four other ASEAN members enter the
priority set. In simulations, EW-DRPL has less than one-third of the regret of the interaction
approach, and equity weighting reduces regret by 36--43 per cent when the planner is strongly
inequality averse.

The paper contributes to three literatures. For the economics of AI and climate change, it moves
from technical potential and single-country carbon studies \citep{liu2022,chen2022,dong2024} to
measured, distribution-specific resilience across a sample dominated by emerging and developing
economies. For climate econometrics, it extends the analysis of heterogeneous climate damages
\citep{burke2015,kalkuhl2020} from the country to the income decile and asks what capability
changes the damages. For the econometrics of treatment choice, it extends equality-minded policy
learning \citep{kitagawa2021,athey2021} to a setting in which the ``treatment'' moderates the effect
of an exogenous shock rather than being randomly assigned itself. It provides an orthogonal score
and a regret bound for that case, together with an implementation in open-source code.

Section~\ref{sec:lit} reviews the related literature. Section~\ref{sec:method} presents the method,
Section~\ref{sec:mc} the Monte Carlo evidence and Section~\ref{sec:data} the data.
Section~\ref{sec:results} reports the empirical results and Section~\ref{sec:discussion} discusses
the implications for policy in emerging and developing economies.
Section~\ref{sec:conclusion} concludes.
